\documentclass[11pt,a4paper]{article}
\usepackage[margin=2.3cm]{geometry}
\usepackage{amsmath,amssymb,amsthm}
\usepackage{booktabs,array}
\usepackage{pgfplots}
\pgfplotsset{compat=1.17}
\usepackage[colorlinks=true,linkcolor=blue!50!black,citecolor=blue!50!black,urlcolor=blue!50!black]{hyperref}
\usepackage{microtype}

\newcommand{\tr}{\operatorname{tr}}
\newcommand{\SL}{\mathrm{SL}(2,\mathbb{C})}

\title{Twisted Kerr-mass vacuum fields: non-integrability of the axial geodesic flow\\[6pt]
\large Numerical results, with computer-assisted certificates near Kerr and over part of the twist range\\
(AI-assisted)}
\author{Sumit Pal\\ Independent researcher \\ ORCID \href{https://orcid.org/0009-0000-9541-2793}{0009-0000-9541-2793}}
\date{DRAFT, 28 November 2026 (not yet externally reviewed)}

\begin{document}
\maketitle

\begin{abstract}
For a one-parameter family of stationary axisymmetric vacuum fields that keep Kerr's mass profile on the
symmetry axis and deform only the twist (parameter $k$, with $k=0$ Kerr and the physical range $0\le k\le 30/7$),
the first-order normal variational equation along the $L_z=0$ timelike geodesic on the lower symmetry axis has a
Zariski-dense monodromy group in $\SL$ for every real $k\in(0,30/7]$, at energy $E=0.8$ and spins
$a\in\{0.3,0.6,0.8,0.9\}$ ($M=1$). This is a \emph{numerical} result: it rests on 50-digit floating-point
argument-principle counts. By the Morales--Ramis--Sim\'o theory it rules out a complete set of first integrals
meromorphic near the complexified orbit, so, at this numerical grade, Carter's constant does not survive the twist.
Three pieces are stronger. (i) Next to Kerr, for $0<|k|\le\varepsilon_a$ with
$\varepsilon_a = 2.79\times10^{-3},\,1.92\times10^{-3},\,8.39\times10^{-4},\,1.06\times10^{-4}$, the statement is
proved, relative to a stated chain of applicability results, by a closed-form $k^2$ coefficient and a certified
remainder bound. (ii) The end sliver $(30/7-\delta_0,30/7]$ is certified in interval arithmetic.
(iii) Computer-assisted interval-arithmetic certificates cover every real $k\in[1/640,3]$ at $a=0.6$ (so, with (i),
all of $0<k\le3$) and $k\in[1/320,3]$ at $a=0.3$. Each node value was confirmed by an independent
50-digit implementation. The work is AI-assisted: the code, analysis and checks were done by AI sessions under
the author's direction, and no external expert has yet reviewed it.
\end{abstract}

\section{Introduction}
Geodesics of the Kerr metric~\cite{Kerr1963} are integrable because of Carter's constant~\cite{Carter1968}, a hidden symmetry not shared by
generic stationary axisymmetric spacetimes. Whether deformations of Kerr keep or lose it matters for the modelling
of extreme-mass-ratio inspirals. Numerical studies of exact ``bumpy'' or non-Kerr spacetimes, such as the Manko--Novikov family~\cite{MankoNovikov1992}, find chaotic orbits and
plateaus in rotation numbers~\cite{GairLiMandel2008,LukesGerakopoulos2010}. Those are statements about sampled
orbits. A different, algebraic route is Morales--Ramis theory~\cite{MoralesRuiz1999,MoralesRamis2001,MRS2007}: if
the variational equation along one particular solution has a differential Galois group whose identity component is
not abelian, then there is no complete set of first integrals meromorphic near that orbit.

That route has been used for static Zipoy--Voorhees fields~\cite{MPS2013}. For the Hartle--Thorne slowly
rotating metric, non-integrability has been shown numerically~\cite{DestounisKokkotas2023}. For undeformed Kerr with a non-spinning orbiting body carrying a tidally induced quadrupole, an analytic proof
shows that no deformation of the geodesic Carter constant stays conserved, for generic tidal couplings and spin,
with numerical evidence of chaos~\cite{Ramond2026}. For equatorially symmetric Newtonian models, Markakis
discusses which multipole structures admit a Carter-like constant~\cite{Markakis2014}.

Here we study a family that keeps Kerr's \emph{mass} profile on the axis exactly and changes only the twist. The
orbit we use runs up and down the lower symmetry axis. On the axis the field is described by Ernst
data~\cite{Ernst1968}, and the normal variational equation is a second-order linear equation in the axis
coordinate. The family is not reflection-symmetric about the equatorial plane, and by black-hole uniqueness no
member with $k\neq0$ is a regular vacuum black hole~\cite{ChruscielCosta2008}.

We state each result at exactly its grade: \emph{proved} (relative to stated inputs), \emph{certified} (by interval
arithmetic, computer-assisted), \emph{numerical} (50-digit floating point, with controls), or \emph{unverified}.

\section{Setting}
\paragraph{Units and coordinates.} $M=1$; the particle is timelike ($\mu=1$), with energy per unit mass $E$. The
coordinate $z$ is the Weyl coordinate shifted by $M$, so that Kerr's upper-axis Ernst potential is
$e_+=(z-ia-2)/(z-ia)$, with real part $f_0=1-2z/(z^2+a^2)$. The rod endpoints are $z=1\pm\sqrt{1-a^2}$, and the
lower outer axis is $z<1-\sqrt{1-a^2}$.

\paragraph{The family.} The upper-axis data keep $f_0$ and replace Kerr's twist $c_0$ by
\begin{equation}
(c_0'')^2=\frac{8\,B(u^2,k)}{5a^6(u^2+1)^6},\qquad u=z/a,\qquad B(w,k)=(90-21k)w^2-(60-18k)w+(10-k),
\end{equation}
with $c_0,c_0'\to0$ as $z\to+\infty$ (asymptotically flat, no NUT charge), and the branch of $\sqrt B$ positive at
large positive $u$. At $k=0$, $B=10(3w-1)^2$ and $c_0$ is Kerr's $-2a/(z^2+a^2)$. The data are real on both outer
axis segments exactly for $0\le k\le 30/7$. $B$ has discriminant $240k(k-4)$, and at $k=30/7$ its leading
coefficient vanishes.

\paragraph{The lower axis and the equation.} The lower-axis data are $e_-=1/e_+$. We derive this for
asymptotically flat data, without reflection symmetry. The lower-axis mass function is $f_0^-=f_0/(f_0^2+c_0^2)$.
Along the $L_z=0$ geodesic on the lower axis, the first-order normal variational equation is
\begin{equation}
v''=-\tfrac{3}{16}\Big(\frac{P'}{P}\Big)^2 v,\qquad P=E^2-f_0^-.
\end{equation}
Its singular points away from the branch structure of $c_0$ are the zeros $r_\pm$ of $P$ and the zeros $p_\pm$ of
$f_0^2+c_0^2$. The equation depends on $c_0$ itself, through $f_0^-$, as well as on $c_0'$.

\paragraph{Loops and the test.} For $k$ up to about 4, $M_0$ is the monodromy around a circle $|z|=R_0$ (with
$R_0=1.1492,1.35,1.3382,1.2639$ at $a=0.3,0.6,0.8,0.9$) that encloses every singularity of $c_0$. $M_1$ is the
monodromy around a capsule lasso that encircles the segment $[r_+,p_+]$. On the circle and the lasso $c_0$ is
single-valued, so both are genuine monodromy elements. Near $k=30/7$ a ``half-plane'' picture is used instead: two
loops inside a disc in $\operatorname{Re}z>0$ that is free of $c_0$'s singularities. With
\[
F(k)=\tr[M_1,M_0]-2,
\]
the conditions $F\neq0$ (irreducible), $\tr M_0\neq0$ and $\tr M_1\neq0$ (not imprimitive), and $\tr M_1$ not the
trace of an element of a finite primitive subgroup (not finite) imply that the Zariski closure of the monodromy
group is $\SL$. By the algebraic-subgroup classification (as in Kovacic~\cite{Kovacic1986}), every proper algebraic
subgroup is reducible, imprimitive or finite.

\section{Results}

\begin{table}[h]
\centering\small
\begin{tabular}{@{}p{5.2cm}p{4.9cm}p{4.6cm}@{}}
\toprule
range of real $k$ (each tested spin) & method & grade\\
\midrule
$(0,30/7-10^{-12}]$ & argument-principle counts over complex $k$-rectangles & numerical (50-digit)\\
$k=30/7$ & direct classification & certified (interval arithmetic)\\
$(30/7-\delta_0,30/7]$ & Zariski-density certificate & certified (interval arithmetic)\\
$(0,\varepsilon_a]$ (Statement A) & closed-form $c_2$ plus a certified remainder bound & proved, relative to the chain of \S3.2\\
$[1/640,3]$ at $a=0.6$; $[1/320,3]$ at $a=0.3$ & segment certificates (Chebyshev and Arb) & certified, computer-assisted\\
$a=0.8$ ($[\varepsilon_a,3]$), $a=0.9$ ($[\varepsilon_a,2]$) & the same & in progress, not reported\\
three degenerate cases & independent loop pairs, located & numerical, post hoc\\
\bottomrule
\end{tabular}
\caption{What is established, and at what grade. $\delta_0=7.4\times10^{-15},\,1.37\times10^{-11},\,
2.4\times10^{-10},\,3.6\times10^{-10}$ at $a=0.3,0.6,0.8,0.9$.}
\label{tab:grades}
\end{table}

\subsection{The numerical result over the whole range}
At each tested spin, the counts show that $F$ has exactly Kerr's double zero at $k=0$ in the near-Kerr rectangle,
and that $F$, $\tr M_0$, $\tr M_1$ and $\tr M_1-s$ for every finite-subgroup trace $s$ are zero-free elsewhere. The
rectangles are $[-1/2,\,\approx4]\times[-1/2,1/2]\,i$ in the circle picture and $[3.95,\,30/7-10^{-12}]\times
[-1/4,1/4]\,i$ in the half-plane picture ($[3.8,\dots]\times[-0.15,0.15]\,i$ at $a=0.9$). The endpoint $k=30/7$ is
classified directly. Three degenerate cases were resolved after the counts, in floating point:
\begin{itemize}
\item at $a=0.8$ the test pair degenerates at $k\approx4.0082$, and an independent pair settles it;
\item at $a=0.9$ an extra zero of $F$ is non-real;
\item at the untested spin $a_c\approx0.858446$ the endpoint pair degenerates, and an independent pair in two
implementations settles it.
\end{itemize}

\subsection{Statement A: next to Kerr, proved relative to the chain}
For $0<|k|\le\varepsilon_a$ the monodromy group is Zariski-dense. The $k^2$ coefficient of $F$ has the closed form
$c_2=\sinh^2(\kappa\omega)\,I_+I_-/\kappa^2$, with $\kappa=a\sqrt{1-E^2}$, built from Kerr's explicit solutions.
Its nonvanishing is certified per spin in interval arithmetic, by a trapezoid rule with the Trefethen--Weideman
error bound~\cite{TrefethenWeideman2014}. We find $|c_2|\approx0.054,\,0.217,\,0.163,\,0.016$. A certified bound
$|F-c_2k^2|<|c_2||k|^2$ on $0<|k|\le\varepsilon_a$ then gives the radius. It comes from an exact $\mathrm{SL}_2$
identity for $F$ and a Dyson bound in Kerr's solution basis, and takes about 80\,s per spin.

The \emph{chain} has three links, and its caveats are part of the statement:
\begin{itemize}
\item the off-axis metric near the orbit (``Proposition 2'') rests on our own majorant-series proof, which is not
yet independently checked;
\item $e_-=1/e_+$ is derived for asymptotically flat data, or equivalently taken as the definition of the lower-axis
data;
\item one cited theorem number in the covering argument is secondhand.
\end{itemize}
Also, $c_2$ changes sign at $a^*\approx0.908$, so Statement A is per spin, not spin-uniform.

\subsection{Interval certificates over part of the twist range}
Each certificate covers a segment $[k_{\rm lo},k_{\rm hi}]$. We interpolate $F$, or $F/k^2$ on dyadic segments near
$k=0$, at $n+1$ Chebyshev points from validated node values computed in Arb~\cite{Arb2017}. A bound $M$ on a
Bernstein ellipse $E_\rho$ ($\rho=4$) is proved by a Dyson bound. ATAP Theorem~8.2~\cite{ATAP} then bounds the
interpolation error by $4M\rho^{-n}/(\rho-1)$, and the conditions follow if the interpolant clears that bound.
Geometric ``holomorphy boxes'' certify that no singular point reaches either loop for any complex $k$ in the
ellipse.

\begin{itemize}
\item \textbf{$a=0.6$.} Every real $k\in[1/640,3]$ is certified, in 35 segments (6 dyadic, 29 of width 0.1), with
no splits. With Statement A, whose interval overlaps it, all of $0<k\le3$ is covered. Production took 15.9\,h on
one core.
\item \textbf{$a=0.3$.} Every real $k\in[1/320,3]$ is certified, in 34 segments. Below $1/320$ the Arb node radii
alone exceed $|F/k^2|$, so the segment $[1/640,1/320]$ cannot be certified at any width. A sliver
$[2.79\times10^{-3},3.125\times10^{-3}]$ above Statement A stays numerical. Production took 14.1\,h.
\item \textbf{$a=0.8$ and $a=0.9$.} Runs over $[\varepsilon_a,3]$ and $[\varepsilon_a,2]$ are in progress, and no
result is reported here. At $a=0.9$ the same floor is reached at $k=1/80$, so a larger gap above Statement A will
remain there.
\item \textbf{Earlier.} The stretches $k\in[1,1.1]$ were certified at $a=0.6$, then independently rechecked, and at
$a=0.9$, then reproduced byte for byte on a second machine.
\end{itemize}
The method stops before the top of the range. Segment centres above $k\approx3.7$--$3.8$ fail, because the
Laurent-series evaluation of $c_0$ used in the circle picture cannot be enclosed tightly enough there. So
$[3,\,30/7-\delta_0]$ remains numerical at every spin.

\begin{figure}[t]
\centering
\begin{tikzpicture}
\begin{loglogaxis}[width=0.92\textwidth,height=7.2cm,xmin=1e-3,xmax=5,ymin=1e-7,ymax=20,
  xlabel={twist parameter $k$},ylabel={certified lower bound on $|F(k)|$},
  legend style={font=\small,at={(0.2,0.97)},anchor=north west},grid=major,grid style={gray!20}]
\fill[green!10] (axis cs:1e-3,1e-7) rectangle (axis cs:1.92e-3,20);
\fill[gray!18] (axis cs:3,1e-7) rectangle (axis cs:4.2857,20);
\addplot[blue,very thick,unbounded coords=jump] table[x=k,y=bound,empty line=jump] {fig_bounds_3_5.dat};
\addlegendentry{$a=0.6$: certified segments}
\addplot[red!70!black,very thick,unbounded coords=jump] table[x=k,y=bound,empty line=jump] {fig_bounds_3_10.dat};
\addlegendentry{$a=0.3$: certified segments}
\draw[dashed,green!50!black] (axis cs:1.92e-3,1e-7) -- (axis cs:1.92e-3,20);
\draw[dashed,red!70!black] (axis cs:2.79e-3,1e-7) -- (axis cs:2.79e-3,20);
\draw[black] (axis cs:4.2857,1e-7) -- (axis cs:4.2857,20);
\node[font=\footnotesize,align=right,anchor=east] at (axis cs:2.9,3e-6) {numerical only $\rightarrow$\\(grey band)};
\node[font=\footnotesize,rotate=90,anchor=south] at (axis cs:4.2857,1e-5) {$k=30/7$};
\node[font=\footnotesize,rotate=90,anchor=center] at (axis cs:1.38e-3,3e-3) {Statement A, $a{=}0.6$};
\end{loglogaxis}
\end{tikzpicture}
\caption{Certified lower bounds on $|F|=|\tr[M_1,M_0]-2|$, one horizontal step per certified segment. On dyadic
segments the certified bound $|F/k^2|\ge b$ is drawn as $|F|\ge b\,k_{\rm lo}^2$. Green: Statement A's interval at
$a=0.6$, $k\le\varepsilon_{0.6}=1.92\times10^{-3}$. Dashed red: $\varepsilon_{0.3}=2.79\times10^{-3}$, below the
first certified $a=0.3$ segment at $1/320$. Grey: $[3,30/7]$, covered numerically (and by the end-sliver
certificate). Near $k=0$, $|F|\approx|c_2|k^2$. Data: \texttt{paper/make\_figure\_data.py}, from the committed
production logs.}
\label{fig:bounds}
\end{figure}

\section{Verification}
\begin{itemize}
\item \textbf{Two implementations.} Every node value of every certified segment was recomputed by an independent
50-digit engine (mpmath) and had to lie inside its Arb enclosure: 36 of 36 gate runs at $a=0.6$, 34 of 34 at
$a=0.3$, and 44 of 44 nodes on the earlier stretches. The two codes share the formula for $c_0''$ and the loop
vertices, and nothing else.
\item \textbf{Must-fail controls,} per mode:
  \begin{itemize}
  \item a planted zero between nodes must be refused;
  \item a held-out node must lie inside the interpolation bound;
  \item an $M$ shrunk by $10^{12}$ must be caught by a sampled-maximum check;
  \item a planted precondition failure must fire.
  \end{itemize}
  For Statement A, the $c_2$ certificate refuses a planted case, and the $\varepsilon$ certifier refuses an
  over-large radius.
\item \textbf{Reproducibility.} Every certified verdict was rerun from its stored node data, and the reruns are
byte-identical. The $[1,1.1]$ certificate at $a=0.9$ was reproduced on a second machine.
\item \textbf{Predictions fixed before the data were examined.} Each run's predictions were committed before it ran,
and hits and misses are recorded.
  \begin{itemize}
  \item At $a=0.3$ all nine predictions hit, including the gap.
  \item At $a=0.6$ every prediction hit except the wall time, which came in under its range.
  \item At $a=0.9$ the predicted lower end of the reach was missed: the floor came at $1/80$, not in
$[7.8\times10^{-4},3.1\times10^{-3}]$.
  \end{itemize}
\item \textbf{Reading.} The interval code was read by an AI session that did not write it, and so were the
derivation of $c_2$ and the $\varepsilon_a$ certificate. Two hostile referee rounds on the argument were run by a
different AI model and answered point by point.
\item \textbf{Loop validity.} A separate floating-point check (50 digits) showed that the circle is a closed loop for the multivalued $c_0$
exactly while it encloses all of $c_0$'s branch points. Past that point it acquires a nonzero elliptic period.
\end{itemize}

\section{Limits}
\begin{itemize}
\item Only the four spins and one energy are tested ($E=1.5$ only on the end sliver at $a=0.6$). The obstruction
weakens roughly like $a^7$--$a^8$ toward low spin, and spins below 0.3 are untested.
\item The excluded integrals are those meromorphic near the complexified orbit. Integrals real-analytic on the
physical exterior that do not continue meromorphically are not excluded.
\item Nothing off the axis is checked: horizons, the rod, singularities and causality are unexamined. These are not
black holes. Null geodesics are untested.
\item The chain of \S3.2 includes our own unrefereed proof, and the numerical part rests on floating-point counts.
\item No professional physicist or mathematician has reviewed this work.
\end{itemize}

\section{Reproducibility}
The software, with the counts, the end-sliver and Statement A outputs, is archived as Ansatz Machine~v1.2.0
(\href{https://doi.org/10.5281/zenodo.24593416}{doi:10.5281/zenodo.24593416}). The segment-certificate production logs and
records are not in that archive yet. They are on the public branch\\
\texttt{personal/2026-11-26-route1-check} of \url{https://github.com/sumit7194/ansatz-machine}.
\begin{itemize}
\item \textbf{Environment.} Python 3.14.5 with mpmath 1.3.0, sympy 1.14.0 and numpy 2.4.6 for the counts. Python
3.13.18 with python-flint 0.9.0 (Arb) for the interval certificates.
\item \textbf{Run times on one laptop core.} Counts take minutes to a few hours, the end-sliver certificates
3--15\,min, the $\varepsilon_a$ certificates about 80\,s, and the segment-certificate productions 14--16\,h, with peaks below
100\,MB.
\item \textbf{Referee guide.} \texttt{paper/REFEREE\_GUIDE.md} maps each claim to its script, expected output and
run time.
\end{itemize}

\section{AI assistance}
This work was carried out with AI assistance. Claude (Anthropic) wrote most of the code, ran the analyses and
drafted the text under the author's direction. Gemini (Google) ran the hostile referee reads. Independent checks
were made by separate AI sessions:
\begin{itemize}
\item reimplementations;
\item a 50-digit engine against the interval code;
\item exact symbolic checks;
\item code reads by a session that did not write the code.
\end{itemize}
These are \emph{not} human peer review, and no external expert has read this work yet. The author chose the
questions, directed and reviewed the work, and takes full responsibility for the content.

\begin{thebibliography}{99}\small
\bibitem{Carter1968} B.~Carter, Global structure of the Kerr family of gravitational fields, Phys.\ Rev.\ 174
(1968) 1559. \href{https://doi.org/10.1103/PhysRev.174.1559}{doi:10.1103/PhysRev.174.1559}
\bibitem{Kerr1963} R.~P.~Kerr, Gravitational field of a spinning mass as an example of algebraically special
metrics, Phys.\ Rev.\ Lett.\ 11 (1963) 237. \href{https://doi.org/10.1103/PhysRevLett.11.237}{doi:10.1103/PhysRevLett.11.237}
\bibitem{Ernst1968} F.~J.~Ernst, New formulation of the axially symmetric gravitational field problem, Phys.\ Rev.\
167 (1968) 1175. \href{https://doi.org/10.1103/PhysRev.167.1175}{doi:10.1103/PhysRev.167.1175}
\bibitem{MoralesRuiz1999} J.~J.~Morales-Ruiz, \emph{Differential Galois Theory and Non-Integrability of Hamiltonian
Systems}, Birkh\"auser (1999). \href{https://doi.org/10.1007/978-3-0348-8718-2}{doi:10.1007/978-3-0348-8718-2}
\bibitem{MoralesRamis2001} J.~J.~Morales-Ruiz and J.-P.~Ramis, Galoisian obstructions to integrability of
Hamiltonian systems, Methods Appl.\ Anal.\ 8 (2001) 33. \href{https://doi.org/10.4310/MAA.2001.v8.n1.a3}{doi:10.4310/MAA.2001.v8.n1.a3}
\bibitem{MRS2007} J.~J.~Morales-Ruiz, J.-P.~Ramis and C.~Sim\'o, Integrability of Hamiltonian systems and
differential Galois groups of higher variational equations, Ann.\ Sci.\ \'Ec.\ Norm.\ Sup\'er.\ 40 (2007) 845.
\href{https://doi.org/10.1016/j.ansens.2007.09.002}{doi:10.1016/j.ansens.2007.09.002}
\bibitem{Kovacic1986} J.~J.~Kovacic, An algorithm for solving second order linear homogeneous differential
equations, J.\ Symb.\ Comput.\ 2 (1986) 3. \href{https://doi.org/10.1016/S0747-7171(86)80010-4}{doi:10.1016/S0747-7171(86)80010-4}
\bibitem{MPS2013} A.~J.~Maciejewski, M.~Przybylska and T.~Stachowiak, Nonexistence of the final first integral in
the Zipoy--Voorhees space-time, Phys.\ Rev.\ D 88 (2013) 064003. \href{https://doi.org/10.1103/PhysRevD.88.064003}{doi:10.1103/PhysRevD.88.064003}
\bibitem{DestounisKokkotas2023} K.~Destounis and K.~D.~Kokkotas, Slowly-rotating compact objects: the
nonintegrability of Hartle--Thorne particle geodesics, Gen.\ Relativ.\ Gravit.\ 55 (2023) 123.
\href{https://doi.org/10.1007/s10714-023-03170-z}{doi:10.1007/s10714-023-03170-z}
\bibitem{Ramond2026} P.~Ramond, Quadrupolar tidal effects destroy the integrability of black hole geodesics:
analytic proof and numerical evidence of chaos, arXiv:2607.27129 (2026).
\bibitem{Markakis2014} C.~Markakis, Constants of motion in stationary axisymmetric gravitational fields, Mon.\ Not.\
R.\ Astron.\ Soc.\ 441 (2014) 2974. \href{https://doi.org/10.1093/mnras/stu715}{doi:10.1093/mnras/stu715}
\bibitem{GairLiMandel2008} J.~R.~Gair, C.~Li and I.~Mandel, Observable properties of orbits in exact bumpy
spacetimes, Phys.\ Rev.\ D 77 (2008) 024035. \href{https://doi.org/10.1103/PhysRevD.77.024035}{doi:10.1103/PhysRevD.77.024035}
\bibitem{LukesGerakopoulos2010} G.~Lukes-Gerakopoulos, T.~A.~Apostolatos and G.~Contopoulos, Observable signature of
a background deviating from the Kerr metric, Phys.\ Rev.\ D 81 (2010) 124005.
\href{https://doi.org/10.1103/PhysRevD.81.124005}{doi:10.1103/PhysRevD.81.124005}
\bibitem{ChruscielCosta2008} P.~T.~Chru\'sciel and J.~L.~Costa, On uniqueness of stationary vacuum black holes,
Ast\'erisque 321 (2008) 195; arXiv:0806.0016.
\bibitem{TrefethenWeideman2014} L.~N.~Trefethen and J.~A.~C.~Weideman, The exponentially convergent trapezoidal rule,
SIAM Rev.\ 56 (2014) 385. \href{https://doi.org/10.1137/130932132}{doi:10.1137/130932132}
\bibitem{ATAP} L.~N.~Trefethen, \emph{Approximation Theory and Approximation Practice}, extended edition, SIAM (2019).
\href{https://doi.org/10.1137/1.9781611975949}{doi:10.1137/1.9781611975949}
\bibitem{Arb2017} F.~Johansson, Arb: efficient arbitrary-precision midpoint-radius interval arithmetic, IEEE Trans.\
Comput.\ 66 (2017) 1281. \href{https://doi.org/10.1109/TC.2017.2690633}{doi:10.1109/TC.2017.2690633}
\bibitem{MankoNovikov1992} V.~S.~Manko and I.~D.~Novikov, Generalizations of the Kerr and Kerr--Newman metrics
possessing an arbitrary set of mass-multipole moments, Class.\ Quantum Grav.\ 9 (1992) 2477.
\href{https://doi.org/10.1088/0264-9381/9/11/013}{doi:10.1088/0264-9381/9/11/013}
\end{thebibliography}
\end{document}
